\documentclass[tate,paper=b5,fontsize=10pt,twocolumn]{jlreq}
\usepackage{lltjext}
\usepackage[most]{tcolorbox}
\begin{document}
\section*{第一問　次の文章を読んで、後の問いに答えよ。}

言葉は、それを使う人の経験と結びついて初めて意味を持つ。辞書に載っている定義は、言葉の意味の\rensuji{ア}\kasen{ごく一部}を示しているにすぎない。

たとえば「海」という語を考えてみよう。港町に育った人にとって、海は毎朝目にする身近な風景であり、潮の匂いや波の音と\rensuji{イ}\kasen{切り離すことができない}。一方、内陸で育った人にとっては、海は写真や映像の中にある遠い存在かもしれない。

このように、同じ言葉であっても、それを受け取る人によって喚起されるものは大きく異なる。私たちが会話の中で誤解を生むのは、多くの場合、\rensuji{ウ}\kasen{このずれ}に気づかないためである。

では、言葉の意味を共有するためには何が必要なのだろうか。それは、相手の経験に対する\rensuji{エ}\kasen{想像力}である。自分の言葉が相手にどのように届くかを考え、相手の言葉の背後にある経験を思い描くこと。その努力なしに、本当の意味での対話は成り立たない。

しかし、想像力にも限界がある。他者の経験を完全に理解することは、原理的に不可能だからである。だからこそ、私たちは言葉を重ね、問いを投げかけ、相手の答えに\rensuji{オ}\kasen{耳を傾け続け}なければならない。

言葉は、それを使う人の経験と結びついて初めて意味を持つ。辞書に載っている定義は、言葉の意味のごく一部を示しているにすぎない。

たとえば「海」という語を考えてみよう。港町に育った人にとって、海は毎朝目にする身近な風景であり、潮の匂いや波の音と切り離すことができない。一方、内陸で育った人にとっては、海は写真や映像の中にある遠い存在かもしれない。

このように、同じ言葉であっても、それを受け取る人によって喚起されるものは大きく異なる。私たちが会話の中で誤解を生むのは、多くの場合、このずれに気づかないためである。

では、言葉の意味を共有するためには何が必要なのだろうか。それは、相手の経験に対する想像力である。自分の言葉が相手にどのように届くかを考え、相手の言葉の背後にある経験を思い描くこと。その努力なしに、本当の意味での対話は成り立たない。

しかし、想像力にも限界がある。他者の経験を完全に理解することは、原理的に不可能だからである。だからこそ、私たちは言葉を重ね、問いを投げかけ、相手の答えに耳を傾け続けなければならない。

言葉は、それを使う人の経験と結びついて初めて意味を持つ。辞書に載っている定義は、言葉の意味のごく一部を示しているにすぎない。

たとえば「海」という語を考えてみよう。港町に育った人にとって、海は毎朝目にする身近な風景であり、潮の匂いや波の音と切り離すことができない。一方、内陸で育った人にとっては、海は写真や映像の中にある遠い存在かもしれない。

このように、同じ言葉であっても、それを受け取る人によって喚起されるものは大きく異なる。私たちが会話の中で誤解を生むのは、多くの場合、このずれに気づかないためである。

では、言葉の意味を共有するためには何が必要なのだろうか。それは、相手の経験に対する想像力である。自分の言葉が相手にどのように届くかを考え、相手の言葉の背後にある経験を思い描くこと。その努力なしに、本当の意味での対話は成り立たない。

しかし、想像力にも限界がある。他者の経験を完全に理解することは、原理的に不可能だからである。だからこそ、私たちは言葉を重ね、問いを投げかけ、相手の答えに耳を傾け続けなければならない。

言葉は、それを使う人の経験と結びついて初めて意味を持つ。辞書に載っている定義は、言葉の意味のごく一部を示しているにすぎない。

たとえば「海」という語を考えてみよう。港町に育った人にとって、海は毎朝目にする身近な風景であり、潮の匂いや波の音と切り離すことができない。一方、内陸で育った人にとっては、海は写真や映像の中にある遠い存在かもしれない。

このように、同じ言葉であっても、それを受け取る人によって喚起されるものは大きく異なる。私たちが会話の中で誤解を生むのは、多くの場合、このずれに気づかないためである。

では、言葉の意味を共有するためには何が必要なのだろうか。それは、相手の経験に対する想像力である。自分の言葉が相手にどのように届くかを考え、相手の言葉の背後にある経験を思い描くこと。その努力なしに、本当の意味での対話は成り立たない。

しかし、想像力にも限界がある。他者の経験を完全に理解することは、原理的に不可能だからである。だからこそ、私たちは言葉を重ね、問いを投げかけ、相手の答えに耳を傾け続けなければならない。

\section*{設問}

\begin{tcolorbox}[colback=white,arc=0pt,boxrule=0.4pt,left=2pt,right=2pt]
問\tatechuyoko{1}　傍線部アについて、筆者がそのように述べる理由を\tatechuyoko{50}字以内で説明せよ。
\end{tcolorbox}

\begin{tcolorbox}[colback=white,arc=0pt,boxrule=0.4pt,left=2pt,right=2pt]
問\tatechuyoko{2}　傍線部イについて、筆者がそのように述べる理由を\tatechuyoko{60}字以内で説明せよ。
\end{tcolorbox}

\begin{tcolorbox}[colback=white,arc=0pt,boxrule=0.4pt,left=2pt,right=2pt]
問\tatechuyoko{3}　傍線部ウについて、筆者がそのように述べる理由を\tatechuyoko{70}字以内で説明せよ。
\end{tcolorbox}

\begin{tcolorbox}[colback=white,arc=0pt,boxrule=0.4pt,left=2pt,right=2pt]
問\tatechuyoko{4}　傍線部エについて、筆者がそのように述べる理由を\tatechuyoko{40}字以内で説明せよ。
\end{tcolorbox}

\begin{tcolorbox}[colback=white,arc=0pt,boxrule=0.4pt,left=2pt,right=2pt]
問\tatechuyoko{5}　傍線部オについて、筆者がそのように述べる理由を\tatechuyoko{50}字以内で説明せよ。
\end{tcolorbox}

\begin{tcolorbox}[colback=white,arc=0pt,boxrule=0.4pt,left=2pt,right=2pt]
問\tatechuyoko{6}　傍線部アについて、筆者がそのように述べる理由を\tatechuyoko{60}字以内で説明せよ。
\end{tcolorbox}

\begin{tcolorbox}[colback=white,arc=0pt,boxrule=0.4pt,left=2pt,right=2pt]
問\tatechuyoko{7}　傍線部イについて、筆者がそのように述べる理由を\tatechuyoko{70}字以内で説明せよ。
\end{tcolorbox}

\begin{tcolorbox}[colback=white,arc=0pt,boxrule=0.4pt,left=2pt,right=2pt]
問\tatechuyoko{8}　傍線部ウについて、筆者がそのように述べる理由を\tatechuyoko{40}字以内で説明せよ。
\end{tcolorbox}

\begin{tcolorbox}[colback=white,arc=0pt,boxrule=0.4pt,left=2pt,right=2pt]
問\tatechuyoko{9}　傍線部エについて、筆者がそのように述べる理由を\tatechuyoko{50}字以内で説明せよ。
\end{tcolorbox}

\begin{tcolorbox}[colback=white,arc=0pt,boxrule=0.4pt,left=2pt,right=2pt]
問\tatechuyoko{10}　傍線部オについて、筆者がそのように述べる理由を\tatechuyoko{60}字以内で説明せよ。
\end{tcolorbox}

\begin{tcolorbox}[colback=white,arc=0pt,boxrule=0.4pt,left=2pt,right=2pt]
問\tatechuyoko{11}　傍線部アについて、筆者がそのように述べる理由を\tatechuyoko{70}字以内で説明せよ。
\end{tcolorbox}

\begin{tcolorbox}[colback=white,arc=0pt,boxrule=0.4pt,left=2pt,right=2pt]
問\tatechuyoko{12}　傍線部イについて、筆者がそのように述べる理由を\tatechuyoko{40}字以内で説明せよ。
\end{tcolorbox}

\begin{tcolorbox}[colback=white,arc=0pt,boxrule=0.4pt,left=2pt,right=2pt]
問\tatechuyoko{13}　傍線部ウについて、筆者がそのように述べる理由を\tatechuyoko{50}字以内で説明せよ。
\end{tcolorbox}

\begin{tcolorbox}[colback=white,arc=0pt,boxrule=0.4pt,left=2pt,right=2pt]
問\tatechuyoko{14}　傍線部エについて、筆者がそのように述べる理由を\tatechuyoko{60}字以内で説明せよ。
\end{tcolorbox}

\begin{tcolorbox}[colback=white,arc=0pt,boxrule=0.4pt,left=2pt,right=2pt]
問\tatechuyoko{15}　傍線部オについて、筆者がそのように述べる理由を\tatechuyoko{70}字以内で説明せよ。
\end{tcolorbox}

\begin{tcolorbox}[colback=white,arc=0pt,boxrule=0.4pt,left=2pt,right=2pt]
問\tatechuyoko{16}　傍線部アについて、筆者がそのように述べる理由を\tatechuyoko{40}字以内で説明せよ。
\end{tcolorbox}

\end{document}
